\begin{frame}{Exemplo: Caso 2 de Hubbers}
    Será considerado aqui o Caso 2 proposto por Hubbers para demonstrar o funcionamento do sistema algébrico computacional e da verificação formal.
    
    \vspace{0.2cm}

    \begin{block}{Aplicação $F$ (Caso 2)}
        \vspace{-0.4cm} 
        \begin{align*}
            F &= (x_1, F_2, x_3, F_4) \\
            F_2 &= x_2 - \frac{1}{3} x_1^3 - h_2 x_1 x_3^2 - q_2 x_3^3 \\
            F_4 &= x_4 - x_1^2 x_3 - h_4 x_1 x_3^2 - q_4 x_3^3
        \end{align*}
    \end{block}
    \vfill
    A ilustração desse processo foi feita a partir de uma animação produzida na \textit{engine} Manim\footnote{\url{https://github.com/3b1b/manim}}.
\end{frame}